% !TeX root = ../papers/nc-p3.tex
\section{正則格子の比較による射影軌道の全正則特殊化}
\label{sec:regular-orbit-valuations}
\begin{theorem}\label{thm:regular-orbit-valuations}
$D$ を離散付値環、$K$ をその分数体とする。二つの標準四窓の $D$ 上の模型を取り、
一般・特殊ファイバーがいずれも標準quadratic Koszul左右AS正則とする。
分数体で一方を他方に運ぶ四文字の基底変換を $g_i$ とし、
標準四位置の文字格子に関する $\operatorname{ord}\det g_i$ を $v_i$ とする。
原始的な最上テンソルの生成元どうしを比較した倍率の付値を $m$ とする。このとき
\[
 \sum_{i=0}^3 v_i=4m,\qquad
 \sum_{i=0}^3\operatorname{ord}\det(g:A_{i,i+n}\longrightarrow A'_{i,i+n})
       =4m\binom{n+3}{4}\quad(n\ge0).
\]
従って、半安定な標準AS窓の射影軌道閉包に属する標準AS窓はすべて半安定で、
半安定商において元の窓と同じ点に写る。とくに多重安定な標準AS窓の
射影軌道閉包に属する標準AS窓は元の添字保持同型類だけである。
\end{theorem}
\begin{proof}[付値のEuler等式]
各模型で窓の二切断の像を二次関係の格子に取る。階数6は一般・特殊で等しく
単元な六次小行列を持つので、これは直和因子である。生成した二次代数の各長さ $n$ の経路加群は
有限表示で、両ファイバーで同じ次元 $h_n=\binom{n+3}3$ だから自由である。
単純の線型解消に使う二次双対の各次数の双対格子 $W_{i,j}$ も、
両ファイバーの次元一致と単元な小行列から自由で、基底変換と可換である。
最高積の左右対合は剰余で完全なので格子上で完全である。
左右Koszul複体の各長さでの項も自由で、特殊ファイバーの完全性から正の長さでは
格子上で分裂完全である（単元なピボットを端から順に除ける）。

長さ $n$ の経路格子の行列式付値を $a_i(n)$ とする。比較は $K$ 上の同型でよく、
行列式は両側の $D$ 基底で取る。$W_{i,2}$ の比較付値を $r_i$ とし、
最上直線 $W_{i,4}$ の付値は原始テンソルの巡回で全位置で $m$ である。
完全対合から $r_i+r_{i+2}=6m$、$W_{i,3}$ の比較付値は $4m-v_{i+3}$ となる。
分裂完全な二列の間の分数体同型では各項の行列式付値の交代和が零である。
各項 $W_{i,j}\otimes A_{i+j,i+n}$ に適用し、始点四つの和を取る。
$A(x)=\sum_{n\ge0}(\sum_i a_i(n))x^n$, $V=\sum_i v_i$ とすれば
\[
 (1-x)^4 A(x)=(1-x)^{-4}\big(Vx-12m x^2+(16m-V)x^3-4m x^4\big).
\]
実際、テンソル積の付値は二因子の付値に相手の階数を掛けた和である。
次数零は同一の単純の比較で寄与しない。

両模型の一次格子の間で $g_i,g_i^{-1}$ の全係数の極の次数を定数 $c$ 以下に取る。
長さ $n$ の積で各経路格子が生成されるので、$g$ は一方を他方の $\pi^{-cn}$ 倍以内に送り、
逆も同様である。従って各基本因子の付値の絶対値は $cn$ 以下、
$|a_i(n)|\le cn h_n=O(n^4)$ である。上式の右の分子は $x=1$ で少なくとも三重に消える。
さもなくば $(1-x)^{-8}$ を掛けた係数の増大がこの評価を超えるためである。
一次微分は $8m-2V$ なので $V=4m$ を得る。分子は $4mx(1-x)^3$ に等しく、
係数を取ると第二式も得る。
\end{proof}
\begin{proof}[軌道閉包]
射影軌道の境界点を特殊点に持つ曲線をその閉包内で取り、正規化・有限体拡大の後、
一般点を基底群の $K$ 点から来るものとできる。
実際、軌道写像の一般ファイバーには有限拡大上の閉点がある。
有限拡大を追加して各四次行列の行列式の四乗根を取れば、スカラー倍で全位置を
特殊線形基底変換に取れる。特殊点に延びる原始的な四テンソルは、
一定な原始テンソルの像を $\pi^m$ で割ったものに単元を掛けて得る。
特殊点も標準正則なら二つの模型に上式が適用でき、$v_i=0$ なので $m=0$ となる。
ゆえに特殊窓には、元のテンソルの特殊線形群の\emph{アフィン}軌道閉包内の非零代表がある。
元が半安定ならある斉次不変式が非零で、その値はこの代表にも保たれる。
全不変式の値が一致するので射影商でも同じ点である。
元のアフィン軌道が閉ならこの代表はその軌道自身に入り、添字保持同型を得る。
\end{proof}

\begin{corollary}\label{cor:isolated-split-component-all-regular}
定理\ref{thm:isolated-split-component}の59次元成分に属する標準Koszul左右AS正則窓は、
不安定境界も含めて調べると六次式の一類で尽きる。
この一類は標準正則軌跡で開かつ閉であり、他の類への特殊化も他の類からの特殊化もない。
\end{corollary}
\begin{proof}
この成分は一つの多重安定な標準窓の射影軌道閉包である。上の定理を直接適用する。
逆向きは同成分の滑らかな開軌道であることから従う。
\end{proof}
